% !TEX root = ../main.tex
\section{Discussion and Recommendations}
\mypara{Networks hosting LDAP}
The public-facing \gls{LDAP} ecosystem shows few signs of centralization. There is a wide range of hosters: the top 10 account for only 24\% of all instances, and 2.2k hosters account for 80\% of \gls{LDAP} servers.
\gls{LDAP} deployment contrasts with other, more centralized ecosystems such as the Web, email, and DNS. Focusing on the scattered hosting of \gls{LDAP}, the variety of issues we found among self-managed solutions indicate that the maintenance of such \gls{LDAP} servers is complex (compared to, \eg a rollover and upgrade mechanism from a cloud provider). In particular, we find much evidence for outdated software.

Our analysis reveals that \gls{LDAP} clusters are present in certain subprefixes across networks of varying sizes. The subprefixes where \gls{LDAP} are found may be shared with different services. From an operator's perspective, it can be advisable to isolate the deployment of \gls{LDAP} servers to a specific network segment, also using firewalls. This practice can enhance security as it allows operators to prevent unauthorized traffic from accessing critical network segments. In a broader sense, different types of customers of a hosting company or \gls{ISP} could be in isolated networks.

We recommend network operators use our tools to regularly survey their attack surface and to increase awareness of servers running on known ports, as also advised by Rapid7~\cite{project_sonar_rapid7}. By doing so, they can identify exposed, vulnerable servers and ultimately initiate campaigns to remove their online presence. Additionally, network operators can use our tools to identify \gls{LDAP} servers and the network segments in which they are operated and isolate them from other critical services or customers. Network operators can also identify where access attempts to \gls{LDAP} servers in their network originate and block unauthorized clients. We recommend that network operators who self-manage solutions perform regular updates or look for solutions with easy upgrade paths.

\mypara{Product lifecycle}
We identified 69.1k \gls{LDAP} products, with 4.3k deprecated. 
The latter servers are (sometimes long) overdue for an upgrade.
Our findings indicate that several \gls{LDAP} servers continue to operate on outdated Windows Server versions, such as 2012 or older, increasing the risk of exposure to known vulnerabilities. Such servers should be migrated to more recent releases, for example, Windows Server 2022, as version 2016 is already under an extended support period.
This also highlights the risk to individuals who rely on organizations that run outdated systems. Users may unknowingly rely on an outdated \gls{LDAP} server where credentials are stored for authentication. Credential theft could occur if such an outdated server is compromised.

Network operators are hence advised to take the task of keeping their networks safe from deprecated services seriously. With our tool set, they can effectively identify deprecated \gls{LDAP} deployments and even extend it to fit their purpose, such as live monitoring or easy access via a dashboard. Our suggestion is to ensure that old \gls{LDAP} software is identified, isolate it in a more specific network segment, and initiate communication with the administrator of the \gls{LDAP} server to patch their software.

\mypara{Communication security}
We analyzed server-side \gls{TLS} support and the certificates used. \gls{LDAP} servers are still using cipher suites that are no longer recommended, and there is a low number of valid certificate chains. We identified a small number of compromised certificates as well. The \gls{TLS} configuration of \gls{LDAP} servers is a metric that an operator can use to determine the security of their network.

Despite the low number of servers using deprecated \gls{TLS} versions, a considerable number of them still do not follow best practices with cipher suites. We recommend that \gls{LDAP} administrators use a robust certificate management system on public-facing \gls{LDAP} servers, such as those for the Web. Use of deprecated \gls{TLS} version, legacy cipher suite, and poor \gls{LDAP} \gls{PKI} are also symptoms of old \gls{LDAP} instances. These metrics help operators to identify unreliable \gls{LDAP} deployments and ultimately enhance the network security by implementing the measures we have recommended so far.

Finally, network operators can also consider going beyond their infrastructure. Ensuring proper management of their own \gls{LDAP} servers contributes to the overall security posture of the Internet, but equally important is awareness of poorly maintained deployments in other networks. For example, organizations that rely on external service providers or federated authentication systems may be indirectly affected by outdated or misconfigured \gls{LDAP} servers outside their administrative domain. Therefore, network operators should be aware that problems in their supply chain may arise.
